\documentclass[conference, a4paper]{IEEEtran}

% =========================================================
% PACKAGES
% =========================================================
\usepackage[T1]{fontenc}
\usepackage{cite}
\usepackage{amsmath,amssymb}
\usepackage{booktabs}
\usepackage{url}
\usepackage{microtype}
\usepackage{listings}
\usepackage[section]{placeins}

% =========================================================
% TITLE
% =========================================================
\title{
Verification-Guided Q-Learning for Web Vulnerability Testing:
An Empirical Evaluation
}

\author{
\IEEEauthorblockN{Kevin Imanuel Setiawan}
\IEEEauthorblockA{
Program Studi Teknologi Informasi, Fakultas Informatika\\
Telkom University Surabaya\\
Surabaya, Indonesia
}
}
\begin{document}

\maketitle

% =========================================================
% ABSTRACT & KEYWORDS
% =========================================================
\begin{abstract}
Automated web vulnerability testing can produce numerous
candidate findings, yet unusual application responses do not
necessarily indicate exploitable weaknesses. This study
evaluates a verification-guided Q-Learning prototype that
selects endpoint--payload actions against OWASP Juice Shop.
The agent observes HTTP responses, derives discrete state
features, and incorporates verification outcomes into reward
assignment. Its configured action space contains 120
endpoint--payload combinations, formed from eight endpoints
and 15 payloads per endpoint. Across five seeds, Q-Learning
verified a mean of 12.8 actions (SD 1.10; range 11--14) and
found all 14 reference actions in one run. A random policy found
all 14 in five runs, with mean full-coverage episode 308.4; a
fixed-order scan found all 14 by episode 66 in one run. A
Q-Learning ablation without verification-guided reward found a
mean of 13.0 actions across five seeds. These comparisons use
the same curated reference and verification procedure, so they
do not establish generalization or independent detector accuracy.
A separate detector benchmark reported recall 0.929 and F1-score
0.963 on 74 labelled actions; 46 unknown cases were excluded.
The results provide no evidence that Q-Learning improves
discovery efficiency in this fixed, source-informed action space.

\end{abstract}
\begin{IEEEkeywords}
Automated penetration testing, reinforcement learning, Q-Learning,
web vulnerability testing, vulnerability verification, OWASP Juice Shop
\end{IEEEkeywords}
% =========================================================
% 1. INTRODUCTION
% =========================================================
\section{Introduction}

Web applications expose many routes and input points, so automated
testing can produce response anomalies that need careful review.
An error, reflected string, or response-size change may identify a
candidate, but does not alone establish exploitability. A confirmed
finding requires evidence suited to the tested vulnerability class.

Reinforcement learning (RL) formulates testing as a sequence of
actions informed by observations and rewards. Tabular Q-Learning
updates values for state--action pairs from observed transitions
\cite{sutton2018}. Prior security work applies RL to grey-box SQL
injection detection, vulnerability exploitation, and autonomous
penetration testing \cite{sqirl2023,almajali2024,li2025}. Reviews
and agent comparisons also show that reported performance depends
on the environment, action space, and evaluation protocol
\cite{rlreview2026,sok2024}.

This study evaluates a verification-guided Q-Learning prototype in
a fixed OWASP Juice Shop test space. Its contributions are:

\begin{itemize}
    \item A tabular Q-Learning workflow that uses verification
    outcomes when assigning rewards.

    \item A comparison with random and fixed-order policies,
    together with reward-verification and reward-scale ablations,
    under a matched target-request budget.

    \item A detector-blind audit using direct response and browser
    evidence, with inconclusive actions retained as unknown.
\end{itemize}

The endpoints, payloads, and surface labels were selected with
source-informed human prior. The findings therefore describe this
configured application and do not establish performance on arbitrary
targets or autonomous endpoint discovery.

% =========================================================
% 2. RELATED WORK
% =========================================================
\section{Related Work}

SQIRL examines grey-box RL for SQL injection detection \cite{sqirl2023};
other studies consider deep RL exploitation, partially observable
agents, and larger network settings \cite{almajali2024,li2025,ghanem2023}.
Surveys and comparative evaluations emphasize that conclusions depend
on each agent's target and test protocol \cite{rlreview2026,sok2024}.
The OWASP Web Security Testing Guide informs the study's distinction
between an anomaly and a finding \cite{owaspwstg}. SQL injection and
file candidates receive controls, while XSS is checked in a browser
with Playwright \cite{playwright}. These checks bound the claims to
the implemented test cases.

\section{Methodology}

\subsection{Experimental Environment}

Trials used OWASP Juice Shop 20.2.0, an intentionally insecure
security-training application \cite{juiceshop}, on Windows with
Node.js v24.22.1. Each fresh instance ran
\texttt{node build/app.js} on port 3001 and reinitialized SQLite.
A counting proxy on port 3000 forwarded target traffic; requests
timed out after 10 seconds. Eight benign probes set endpoint-specific
response-length baselines.

The agent selects from eight configured routes and 15 payloads
per route, observes each HTTP response, and updates state--action
values. Results are specific to this configuration.

\subsection{HTTP API and Endpoint Coverage}

The client places each payload in a configured query parameter,
path, or JSON body. Status, body length, content type, and response
time form the observation. Table~\ref{tab:api-endpoints} lists the
eight routes.

\begin{table}[!t]
\caption{HTTP Routes Used in the Experiment}
\label{tab:api-endpoints}
\centering
\footnotesize
\setlength{\tabcolsep}{3pt}
\begin{tabular}{c p{0.41\columnwidth} p{0.43\columnwidth}}
\toprule
\textbf{Idx.} & \textbf{Method and route} & \textbf{Role in the experiment}\\
\midrule
0 & GET \url{/rest/products/search?q=} & Product search; SQLi and search-page XSS checks\\
1 & GET \url{/ftp/} + path & File-serving route; path and extension-bypass checks\\
2 & GET \url{/rest/products/1/reviews?id=} & Review retrieval; plain-surface control\\
3 & GET \url{/rest/track-order/} + path & Order tracking; browser-based XSS checks on the result page\\
4 & POST \url{/rest/user/login} & Login request with a JSON body; SQLi checks\\
5 & PUT \url{/rest/products/1/reviews} & Review submission with a JSON body; plain-surface control\\
6 & POST \url{/api/Feedbacks} & Feedback JSON body; requires CAPTCHA fields\\
7 & GET \url{/api/Challenges?q=} & Challenge API; plain-surface control\\
\bottomrule
\end{tabular}
\end{table}
\setlength{\tabcolsep}{6pt}

Source inspection informed route selection: search and login reach
SQL handlers, \texttt{/ftp/} serves files, and the track-order
identifier is rendered in a browser. Other routes are controls. This
is a researcher-defined surface, not autonomous endpoint discovery.
Feedback requests include a CAPTCHA answer from \url{/rest/captcha}.
Before testing, \texttt{hello} establishes each baseline. XSS checks
load the corresponding single-page application (SPA). Source-informed
surface labels select checks; they do not predetermine findings.

\subsection{Action Space}

Let $\mathcal{E}$ and $\mathcal{P}$ be the configured endpoint
and payload sets. Their Cartesian product gives
$\mathcal{A}=\mathcal{E}\times\mathcal{P}$ and
$|\mathcal{A}|=8\times15=120$ actions.

Table~\ref{tab:payloads} lists the templates, which are paired with
each route and placed in its configured input. The long-input probe
contains 1,628 ASCII \texttt{A} characters. This action space is fixed
before training.

\begin{table*}[!t]
\caption{Configured Payload Templates}
\label{tab:payloads}
\centering
\scriptsize
\begin{tabular}{p{0.20\textwidth}p{0.12\textwidth}p{0.63\textwidth}}
\toprule
\textbf{Name} & \textbf{Class} & \textbf{Payload}\\
\midrule
\texttt{sqli\_tautology} & SQLi & \lstinline|' OR '1'='1|\\
\texttt{sqli\_comment} & SQLi & \lstinline|apple'--|\\
\texttt{sqli\_union} & SQLi & \lstinline|' UNION SELECT NULL--|\\
\texttt{sqli\_time} & SQLi & \lstinline|' OR SLEEP(5)--|\\
\texttt{sqli\_wildcard} & SQLi & \lstinline|%s%s%s%s%s%s%s%s|\\
\texttt{sqli\_login} & SQLi & \lstinline|{"email":"admin@juice-sh.op' OR '1'='1","password":"a"}|\\
\texttt{xss\_img} & XSS & \lstinline|<img src=x onerror=alert(document.domain)>|\\
\texttt{xss\_script} & XSS & \lstinline|<script>alert(document.domain)</script>|\\
\texttt{xss\_svg} & XSS & \lstinline|<svg onload=alert(1)>|\\
\texttt{xss\_review} & XSS & \lstinline|{"message":"<img src=x onerror=alert(document.domain)>","author":"rlpoc"}|\\
\texttt{lfi\_poison\_null} & File & \lstinline|package.json.bak%2500.md|\\
\texttt{lfi\_coupons} & File & \lstinline|coupons_2013.md.bak%2500.md|\\
\texttt{lfi\_wap\_rules} & File & \lstinline|suspicious_errors.yml%2500.pdf|\\
\texttt{fuzz\_long} & Fuzzing & \texttt{A} repeated 1,628 times\\
\texttt{fuzz\_null} & Fuzzing & \lstinline|admin%00|\\
\bottomrule
\end{tabular}
\end{table*}

Each action submits one template to one route. The following state
is derived from its response. Payload and surface choices encode
prior knowledge of Juice Shop and constrain the experiment's scope.

\subsection{State Representation}

The state contains eight features:

\begin{equation}
\begin{split}
s_t = (&\text{status\_cat},
\text{length\_cat},
\text{err\_flag},
\text{refl\_flag},\\
&\text{vuln\_ind},
\text{ep\_idx},
\text{resp\_time\_cat},
\text{content\_cat}).
\end{split}
\label{eq:state}
\end{equation}

The features encode HTTP status and length categories, error and
reflection flags, a vulnerability indicator, endpoint index, response
time, and content type. Length changes are measured against the route's
baseline; changes above 15\% are significant. Response time is fast
below 0.5 s, medium from 0.5--2.0 s, and slow above 2.0 s. The state
captures only these configured response characteristics.

\subsection{Q-Learning Update}

The tabular update uses the observed reward and the largest estimated
value in the next state \cite{sutton2018}:

\begin{equation}
\begin{aligned}
Q(s_t,a_t) \leftarrow Q(s_t,a_t) + \alpha[&r_t + \gamma \max_{a'} Q(s_{t+1},a')\\
&- Q(s_t,a_t)].
\end{aligned}
\label{eq:ql}
\end{equation}

Here $s_t,a_t,r_t,s_{t+1}$ denote the current state, selected action,
reward, and next state; $\alpha$ and $\gamma$ are the learning rate
and discount factor.

The configuration uses $\alpha=0.1$, $\gamma=0.9$, and an
$\epsilon$-greedy policy with $\epsilon$ decaying from 1.0 by 0.90
per episode to a 0.1 floor. This floor is reached at episode 23; only
the first seven episodes have $\epsilon\geq0.5$. The decay was not
tuned. Only verified findings are masked; other actions remain
eligible. Episode counts vary under the shared request cap because
verification traffic differs.

\subsection{Comparative Experimental Protocol}

Each trial had an exact cap of 1,640 requests forwarded to Juice Shop,
counting probes, baselines, CAPTCHA lookups, controls, browser
navigations, and loaded resources. Excess attempts were blocked and
logged. Q-Learning and random selection used five seeds (42--46);
fixed order was repeated on five fresh instances. It cycled through
ascending action indices, skipping only verified actions. All policies
used the same cap and mask.

The five-seed ablation kept verification for reporting and masking,
but used the pre-verification reward for Q updates. Reward sensitivity
used three paired seeds (42--44) with verified/anomaly/normal/error
rewards of $+25/+5/-1/-5$ and $+100/+20/-1/-5$; length bonuses were
unchanged. Each trial began from a fresh instance. Actions are
conditioned on the preceding response summary, not on application
workflow state or an exploit chain; the method therefore tests
response-informed scheduling within this fixed action set.

\subsection{Verification-Guided Reward Assignment}

Rewards are $+50$ for verified findings, $+10$ for unverified
anomalies, $-1$ for ordinary responses, and $-5$ for failed attack
requests, plus a length-category bonus ($+0$, $+0$, $+5$, or $+10$).
SQLi and file candidates must differ from both a benign \texttt{hello}
control and a distinct neutralized-input control. For XSS, headless
Chromium checks whether an instrumented DOM canary executes while the
neutralized control remains inactive; reflection is not enough. A
failed control leaves the candidate inconclusive and unverified.
After demotion, the base reward is $+10$ only if anomaly flags persist;
otherwise it is $-1$, with the length bonus retained.

The analyzer's SQLi candidates require a database error with an HTTP
error, or a successful response-length change of at least 15\% for a
SQL-like input; repeated wildcards have a separate rule. File disclosure
requires protected-content evidence or a \texttt{\%2500} extension
bypass with a disallowed content type; a listing or size change alone
does not qualify. The probes target three named backup files. Training
logs lack raw bodies and content types, so the audit collected these
evidence items separately. Reward values were configured rather than
optimized; two alternative scales are compared in
Section~\ref{sec:reward-sensitivity}.

\subsection{Detector-Blind Oracle Audit}

The authors audited all 120 actions without importing the response
analyzer or prior labels. Evidence included returned records, SQLite
errors, authentication outcomes, byte-level protected-file matches,
and five browser-dialog attempts with neutralized controls. The
version-specific Juice Shop challenge catalogue corroborated, but did
not determine, labels~\cite{juicechallenges}. The audit found 13
positive, 69 negative, and 38 unknown actions; inconclusive cases remain
unknown. Raw bodies, hashes, content types, DOM snapshots, screenshots,
and per-attempt records are retained. The procedure was detector-blind
but author-run; no external reviewer adjudicated the labels.

\subsection{Evaluation Metrics}

Comparisons use true positives (TP), false positives (FP), false
negatives (FN), and true negatives (TN). Recall and F1-score are
computed on labelled actions only; the 38 unknown actions are excluded.

\begin{equation}
\mathrm{Recall} = \frac{TP}{TP+FN}.
\label{eq:recall}
\end{equation}

\begin{equation}
\mathrm{F1} =
2\frac{\mathrm{Precision}\times\mathrm{Recall}}
{\mathrm{Precision}+\mathrm{Recall}}.
\label{eq:f1}
\end{equation}

% =========================================================
% 4. RESULTS AND DISCUSSION
% =========================================================
\section{Results and Discussion}

\subsection{Comparative Results}

Table~\ref{tab:strategy-results} summarizes all 22 trials.
Each trial used 1,585 primary endpoint--payload episodes.
The 14-action reference set is the same curated set used in the
detector benchmark; it is not an independent oracle. Q-Learning
found a mean of 12.8 actions (SD 1.10; range 11--14), and only
one of five seeds found all 14. Uniform random selection found
all 14 in all five seeds, while the fixed-order scan found all
14 by episode 66 in its single deterministic run. The random
baseline's mean episode for full coverage was 308.4 (range
217--440). These results provide no evidence that Q-Learning
improves discovery efficiency in this fixed action space.

\begin{table*}[!t]
\caption{Matched Primary-Action Budget Comparison}
\label{tab:strategy-results}
\centering
\scriptsize
\begin{tabular}{p{0.20\textwidth}c p{0.28\textwidth}c c p{0.17\textwidth}}
\toprule
\textbf{Method} & \textbf{$n$} & \textbf{Verified actions, mean (SD); range} & \textbf{All 14} & \textbf{Episode to all 14$^*$} & \textbf{Mean requests}\\
\midrule
Q-Learning & 5 & 12.80 (1.10); 11--14 & 1/5 & 787.0 & 1,634.2\\
Random & 5 & 14.00 (0.00); 14--14 & 5/5 & 308.4 & 1,639.2\\
Fixed order & 1 & 14.00; 14--14 & 1/1 & 66.0 & 1,639.0\\
QL without verification reward & 5 & 13.00 (1.22); 11--14 & 2/5 & 989.0 & 1,635.6\\
Lower reward & 3 & 13.33 (0.58); 13--14 & 1/3 & 787.0 & 1,637.0\\
Higher reward & 3 & 13.33 (0.58); 13--14 & 1/3 & 787.0 & 1,636.7\\
\bottomrule
\end{tabular}
\end{table*}

The fixed-order scan visits action indices in ascending order;
the final positive reference action occurs at index 65, so all 14
were reached by episode 66. Its single run is descriptive rather
than a variance estimate. The random policy reached full coverage
in all five seeds. Its performance is a direct baseline for the
configured 120-action space, not evidence about unknown endpoints
or vulnerabilities outside that space.

The Q-Learning-without-verification-reward ablation averaged
13.0 verified actions (SD 1.22; range 11--14), with full coverage
in two of five seeds. This ablation still used verification for
the finding log and action mask; it removes the verification
verdict from the reward update, not from the whole policy.

The 14 reference actions are five search-SQLi, two search-XSS,
three FTP file-disclosure (labelled fuzzing by the detector),
three track-order-XSS, and one login-SQLi action. These are
repeated endpoint--payload findings across five endpoint--class
pairs, not 14 independent root causes.
\subsection{Reward Sensitivity}
\label{sec:reward-sensitivity}

The lower ($+25/+5/-1/-5$) and higher ($+100/+20/-1/-5$)
reward settings produced the same verified-action counts as the
base reward ($+50/+10/-1/-5$) for paired seeds 42--44: 14, 13,
and 13 actions, respectively. In each setting, only seed 42
found all 14, at episode 787. This three-seed sensitivity check
shows no observed change in these discovery outcomes under the
tested reward scales; it does not establish insensitivity for
other seeds, reward ratios, or applications. Length bonuses
were unchanged.

\subsection{Detector-Blind Audit Results}

The detector-blind audit classified 13 actions as positive, 69 as
negative, and 38 as unknown. Previously recorded no-agent analyzer
outputs, rescored against these labels, matched all 13 positives and all 69
negatives; the 38 unknown cells were excluded. This subset agreement
is not an external accuracy estimate: the audit was author-run, and
the RL verifier showed a reproducibility discrepancy for one SVG
action, discussed below.

\begin{table*}[!t]
\caption{Direct-Effect Evidence and Negative Controls}
\label{tab:oracle-evidence}
\centering
\scriptsize
\begin{tabular}{p{0.16\textwidth}p{0.45\textwidth}p{0.31\textwidth}}
\toprule
\textbf{Test group} & \textbf{Observed response or browser effect} & \textbf{Control result}\\
\midrule
Search SQLi (5 actions) & Tautology returned all 46 catalogue records; three crafted inputs returned SQLite parser errors; wildcard input returned 16 records. & Benign search returned 0; wildcard-neutralized search returned 0; neutralized SQL inputs produced no parser error.\\
Login SQLi (1 action) & Response contained an admin authentication token (redacted in the saved record). & Invalid-password control returned HTTP 401 and no token.\\
File disclosure (3 actions) & Raw bodies matched the packaged backup, coupon, and YAML files by SHA-256; sizes were 4,440, 142, and 741 bytes, with octet-stream or YAML content types. & Removing the encoded null-byte bypass produced HTTP 404 for all three requests.\\
XSS (4 actions) & Four original vectors across search and track-order pages triggered a JavaScript dialog in 5/5 attempts each. & Each neutralized control triggered 0/5 dialogs. The separate track-order SVG vector triggered 0/5 and remained negative.\\
\bottomrule
\end{tabular}
\end{table*}

The raw HTTP bodies and content types, hashes, controls, browser DOM
snapshots, screenshots, and five-attempt traces are stored with the
revision evidence. The action-level audit counts the track-order SVG
as negative because the original vector produced no dialog in five
direct browser attempts. In several RL runs, the instrumented canary
accepted that action. This difference shows that changing a handler
to a canary can yield a verdict that the original vector does not
reproduce; those policy detections are reported as false positives
against the direct-effect audit.

\subsection{Limitations}

The study covers one Juice Shop release and a human-defined set of
eight routes and 15 payloads. Its state records selected HTTP features,
not application workflow; results do not establish transfer to other
targets. The detector-blind oracle was run by the authors without an
external second reviewer, and 38 cells remain unknown. Any metric is
therefore a consistency measure on this labelled subset.

Stochastic policies use five seeds and reward-scale variants use three
paired seeds. Fixed order is deterministic, repeated on five fresh
instances, and is descriptive rather than a source of stochastic
variance. The shared budget fixes forwarded HTTP requests, so primary
action counts differ; requests beyond the cap are blocked locally.
The launch details for the earlier single-seed run were not recorded;
the replicated trials use the setup stated above.

The state update and verification interfaces could be adapted after
defining a new target's routes, input locations, payloads, baselines,
and verification rules. The current route map and browser pages are
Juice Shop specific, so the prototype is not a general-purpose scanner.

\section{Conclusion}

This study compared verification-guided Q-Learning with
random and fixed-order action selection on a 120-action
OWASP Juice Shop test space. Across five seeds, Q-Learning
verified a mean of 12.8 of the 14 curated reference actions
(SD 1.10; range 11--14), and reached all 14 in one run.
Random selection reached all 14 in five of five runs, with
mean full-coverage episode 308.4. The fixed-order scan reached
all 14 by episode 66 in its single run. The Q-Learning ablation
without verification-guided reward averaged 13.0 actions across
five seeds. These results do not show a discovery-efficiency
advantage for Q-Learning in this finite, source-informed space.

The separate detector benchmark reported recall of 0.929
and F1-score of 0.963 on the labelled subset, with 46 unknown
actions excluded. Because the manual reference checks use
evidence related to the prototype's verification procedure,
these values show consistency with the curated reference set
rather than independent generalization. The comparison uses the
same verifier and labels for all policies, so it does not provide
independent confirmation of vulnerability truth.

Future work should obtain an independent label audit, increase
the number of seeds, enforce an exactly matched total-request
budget including verification overhead, capture raw response and
browser traces, and evaluate additional applications after their
endpoint maps and verification rules have been specified.

% =========================================================
% ACKNOWLEDGMENT
% =========================================================
\begin{thebibliography}{99}

\bibitem{sqirl2023}
A. Al Wahaibi, S. Foley, and S. Maffeis,
``SQIRL: Grey-Box Detection of SQL Injection Vulnerabilities
Using Reinforcement Learning,''
in \textit{Proceedings of the 32nd USENIX Security Symposium},
2023, pp. 6097--6114.

\bibitem{rlreview2026}
J. Liu et al.,
``Autonomous penetration testing using reinforcement learning:
A review and perspectives,''
\textit{Expert Systems with Applications},
vol. 300, Art. no. 130219, 2026,
doi: 10.1016/j.eswa.2025.130219.

\bibitem{sok2024}
M. Simon and S. Mees,
``SoK: A Comparison of Autonomous Penetration Testing Agents,''
in \textit{Proceedings of the 19th International Conference on
Availability, Reliability and Security (ARES)}, 2024,
doi: 10.1145/3664476.3664484.

\bibitem{sutton2018}
R. S. Sutton and A. G. Barto,
\textit{Reinforcement Learning: An Introduction},
2nd ed. Cambridge, MA, USA: MIT Press, 2018.

\bibitem{juiceshop}
OWASP Foundation,
``OWASP Juice Shop.''
[Online]. Available:
\url{https://owasp.org/www-project-juice-shop/}

\bibitem{juicechallenges}
OWASP Foundation,
``OWASP Juice Shop challenge catalogue, version 20.2.0.''
[Online]. Available:
\url{https://github.com/juice-shop/juice-shop/blob/v20.2.0/data/static/challenges.yml}

\bibitem{owaspwstg}
OWASP Foundation,
``Web Security Testing Guide.''
[Online]. Available:
\url{https://owasp.org/www-project-web-security-testing-guide/}

\bibitem{playwright}
Microsoft,
``Playwright Documentation.''
[Online]. Available:
\url{https://playwright.dev/docs/intro}

\bibitem{almajali2024}
A. AlMajali, L. Al-Abed, K. M. A. Yousef,
B. J. Mohd, Z. Samamah, and A. Abu Shhadeh,
``Automated Vulnerability Exploitation Using Deep Reinforcement
Learning,''
\textit{Applied Sciences},
vol. 14, no. 20, Art. no. 9331, 2024,
doi: 10.3390/app14209331.

\bibitem{li2025}
Z. Li et al.,
``EPPTA: Efficient partially observable reinforcement learning
agent for penetration testing applications,''
\textit{Engineering Reports},
vol. 7, no. 1, Art. no. e12818, 2025,
doi: 10.1002/eng2.12818.

\bibitem{ghanem2023}
M. C. Ghanem, T. M. Chen, and E. G. Nepomuceno,
``Hierarchical reinforcement learning for efficient and effective
automated penetration testing of large networks,''
\textit{Journal of Intelligent Information Systems},
vol. 60, pp. 281--303, 2023,
doi: 10.1007/s10844-022-00738-0.

\end{thebibliography}

\end{document}
